% Issue 34: the selected-scale branch receipts; integration is owned elsewhere.
\section{Branch moments and trace washout}
\label{sec:ex-trace-branch-moments}

Fix a closed subgroup of a finite-dimensional unitary group whose real Lie
algebra is nonabelian and simple; the group itself may be disconnected.
Work on one selected history and with the same three tensor families and
chart as in the physical stress construction.
All spaces below are finite dimensional. The four coordinate axes form the
set $I=\{0,1,2,3\}$.
Let $E_4$ be the complex space of grade-four labels, let $\mathcal F_4^{\mathrm{int}}$
be the grade-four field space of the selected interacting table, and let
$\mathcal F_4^0$ be that of the standard table. The maps are complex linear
unless stated otherwise. Write $\pi:E_4\to\mathcal F_4^{\mathrm{int}}$ and
$\pi_0:E_4\to\mathcal F_4^0$ for the projections, $\mathfrak c$ and
$\mathfrak c_0$ for centering, and $\mathbf 1_4,\mathbf1_4^0$ for the
corresponding unit fields. Thus
$\mathfrak c f=f-\langle f\rangle\mathbf1_4$. The centering map is
idempotent and its range is the space of centered fields, using the table
unit law $\langle\mathbf1_4\rangle=1$; the standard table has the same
unit law.

The branch set is $\{0,-,+\}$: $0$ is the principal branch, $-$ is the
false correction and $+$ is the true correction. In the underlying
finite-volume construction these are respectively the selected principal
field and the two fixed correction fields. To specify the choice, let
$J_{ac},Q_d,H_{acd}$ be the selected tensor, trace-current and antisymmetry
fields at a finite volume. Let a superscript $[r]$ denote the same finite
localization of each field and let $\nabla_d$ be the forward lattice
difference. At localization radius $r$ the principal field is
\[
 J^0_{ac}=\begin{cases}
 \sum_d\nabla_d Q_d^{[r]}-\sum_{j\ne0}J_{jj}^{[r]},&a=c=0,\\
 J_{ca}^{[r]}-\sum_d\nabla_d H_{cad}^{[r]},&c<a,\\
 J_{ac}^{[r]},&\text{otherwise}.
 \end{cases}
\]
At scale $k$ the principal field's common localization radius is
$\varrho_k=1000\lceil(\log\widehat\beta_k)^2\rceil$, where $\widehat\beta_k$
is the positive clock introduced below; here $r=\varrho_k$. In particular $\sum_aJ^0_{aa}=\sum_d\nabla_dQ_d^{[r]}$ exactly. For the two
corrections, let $B_{ad,cd}$ be the real inner product of the two lattice
curvatures with the common index $d$, evaluated at the same lattice site.
Their components are
\[
 J^-_{ac}=\begin{cases}0,&a=c,\\\sum_d B_{ad,cd},&a\ne c,\end{cases}
 \qquad
 J^+_{ac}=\begin{cases}\frac14\sum_{j,d}B_{jd,jd},&a=c,\\
 0,&a\ne c.\end{cases}
\]
These formulas are translated from the origin to the relevant lattice site.
The three branches use the same extraction ultrafilter, finite-volume
choices and scale history, with no new subsequence; their recorded labels may
differ. At scale $k$, let $x^b_k(a,c)\in E_4$ be the recorded
label of $J^b_{ac}$, with $a,c\in I$. Let $\mathsf R$ be the
real-linear finite signed-symmetry average on arrays of labels and set
\[
 \operatorname{sym}(x)_{ac}=\tfrac12(x_{ac}+x_{ca}),\qquad
 t^b_k=\mathfrak c\,\pi\mathsf R\operatorname{sym}(x^b_k).
\]
The notation means componentwise application of $\pi$ and $\mathfrak c$.
The symmetry average permutes the axes and applies the corresponding label
map, with the two tensor signs; it includes the real-structure average.
Its normalized weights sum to one. The table projection intertwines these
label motions with the corresponding real-linear field motions, and those
field motions are fixed, bounded maps. In particular, for a diagonal trace,
reindexing the axes cancels the squared signs. Put
$\operatorname{tr}t=\sum_{a\in I}t_{aa}$.

For a triple $w=(w_1,w_2,w_3)\in\mathbb R^3$, define the label-valued
material tensor by
\begin{equation}
 \mathcal S(w)_{ac}=\begin{cases}
 w_1\sum_j P_{aj,aj}+w_2\sum_{i,j}P_{ij,ij},&a=c,\\
 w_3\sum_j P_{aj,cj},&a\ne c,
 \end{cases}
 \qquad
 M(w)_{ac}=\pi_0\mathcal S(w)_{ac}.
 \label{eq:ex-branch-material}
\end{equation}
Here $P_{ij,kl}\in E_4$ is the coordinate representative of the label
whose value on a Taylor jet is the real color inner product
$\langle F_{ij},F_{kl}\rangle$; no additional factor is included. The
axis sums are ordered. There is no
scalar-unit term in \eqref{eq:ex-branch-material}; the full four-component
coefficient is $(0,w_1,w_2,w_3)$. The map $w\mapsto\mathcal S(w)$ is
injective. To see this, let $v$ be a nonzero real color vector. On a jet
with only $F_{01}=v$ and $F_{10}=-v$, the $(2,2)$ and $(0,0)$ entries
are respectively $2w_2\langle v,v\rangle$ and
$(w_1+2w_2)\langle v,v\rangle$; if the tensor is zero, $w_2=w_1=0$.
On a jet with $F_{02}=F_{12}=v$ and their antisymmetric partners, its
$(0,1)$ entry is $w_3\langle v,v\rangle$, so $w_3=0$. Exchange of the
curvature pairs shows $\mathcal S(w)_{ac}=\mathcal S(w)_{ca}$. A signed
permutation sends $F_{ij}$ to $\epsilon_i\epsilon_jF_{\sigma(i)\sigma(j)}$.
In a contracted pair the repeated sign squares to one, and reindexing the
sums gives the two tensor signs. The diagonal terms transform likewise, and
both the curvature-pair labels and real coefficients are fixed by the real structure. Thus the signed and real
symmetry average fixes $\mathcal S(w)$. The branch coefficients are
\begin{equation}
 w^0=(-1,\tfrac14,-1),\qquad
 w^-=(0,0,1),\qquad w^+=(0,\tfrac14,0).
 \label{eq:ex-branch-coefficients}
\end{equation}
Define the normalized scalar field
$\mathcal O=\pi_0(\tfrac14\sum_{i,j}P_{ij,ij})$. Taking the four
diagonal entries in \eqref{eq:ex-branch-material} gives, for every $w$,
\begin{equation}
 \operatorname{tr} M(w)=4(w_1+4w_2)\mathcal O.
 \label{eq:ex-branch-material-trace}
\end{equation}
Consequently
\begin{equation}
 \operatorname{tr}M(w^0)=\operatorname{tr}M(w^-)=0,\qquad
 \operatorname{tr}M(w^+)=4\mathcal O.
 \label{eq:ex-branch-three-traces}
\end{equation}
We use the standard-table one-point fact that each curvature-pair label
has expectation zero and the unit has expectation one. In the free table
the pair labels are Wick ordered. Thus the expectation of each
$M(w)_{ac}$ is zero because its unit coefficient is zero.
Hence $\mathfrak c_0\pi_0\mathcal S(w)=M(w)$.

Let $k_i$ be the selected level, tending to infinity with the chart index
$i$, and let $\widehat\beta_k>0$ and $a_k>0$ be its inverse-coupling clock
and mesh. Denote by $B_k:E_4\to E_4$ the selected transport, by $U_i:\mathcal F_4^0\to\mathcal F_4^{\mathrm{int}}$
the chart isomorphism, and by $\omega_i>0$ its grade-four normalization factor. Explicitly
$\omega_i=\beta_Z(k_i)r_i^4$, where $r_i$ is the chart radius and $\beta_Z$ is
the associated effective coupling. More explicitly, if $A_k$ is the selected label
automorphism of $E_4$, $K$ is the core projection and $\vartheta_k>0$ is
the threshold, then
\[
 B_k=\vartheta_k^{-2}K A_k,\qquad
 V_{k,c}=\pi\, (c\vartheta_k^2 A_k^{-1}).
\]
The centered moment map and the scaled field map are
\begin{equation}
 \mathsf M_i=U_i^{-1}\mathfrak c V_{k_i,\omega_i}:E_4\to\mathcal F_4^0,
 \qquad \mathsf P_i=\omega_iU_i^{-1}:\mathcal F_4^{\mathrm{int}}\to\mathcal F_4^0.
 \label{eq:ex-branch-moment-map}
\end{equation}
The transport and chart normalization give, for all sufficiently large $i$,
\begin{equation}
 \mathsf M_i B_{k_i}x=\mathsf P_i\mathfrak c\pi x
 \quad(x\in E_4).
 \label{eq:ex-branch-chart-identity}
\end{equation}
To see the normalization, $V_{k,c}B_kx=c\pi A_k^{-1}KA_kx$.
The difference between $A_k^{-1}KA_kx$ and $x$ is a multiple of
$A_k^{-1}e_0$, where $e_0$ is the scalar coordinate label. At the selected
cutoff $A_k$ sends the reference unit label to $a_k^4e_0$. The reference
unit projects to $\mathbf1_4$, so $\pi A_k^{-1}e_0=a_k^{-4}\mathbf1_4$
and centering removes the difference. Applying $U_i^{-1}$ yields the identity.
The cutoff after which this holds is fixed before taking $i\to\infty$.
Let $W^b_k=\mathsf R(B_kx^b_k)$, componentwise, and
$W^b=\mathcal S(w^b)$.
The following are the incoming chart limits, stated in their actual
limit order:
\begin{align}
 W^b_k&\longrightarrow W^b &&(k\to\infty,\ b=0,-,+),
 \label{eq:ex-branch-coefficient-input}\\
 \mathsf M_i&\longrightarrow \mathsf M_\infty:=\mathfrak c_0\pi_0
 &&\text{in operator norm as }i\to\infty.
 \label{eq:ex-branch-moment-input}
\end{align}
These limits depend on the selected finite-volume transport and chart
estimates. In the first, the separation-ultrafilter limit is taken at fixed
$k$ before $k\to\infty$.
The existing weighted-estimate construction uses the same reference
coefficient limit, and its quantitative inputs remain subject to the gaps
recorded there. We record the actual finite-volume passage yielding
\eqref{eq:ex-branch-coefficient-input}. Let $T:\mathbb R^3\to
E_4^{I\times I}$ send $w$ to $\mathcal S(w)$. For each $k$ let
$R_k(u)$ be the recorded finite-volume coefficient array before symmetry
averaging and $C^b_k(u)$ the corresponding finite-volume branch
coefficient, with $u$ the separation parameter. Equip these arrays with
the finite-product supremum norm induced by the coordinate norm.
The incoming estimates say that the nonnegative loss
\[
 \delta_k=\exp\!\left[-\tfrac14(\log\log\widehat\beta_k)^2\right]
\]
satisfies, eventually in $k$ and on an ultrafilter-large set of $u$,
\[
 \|R_k(u)-C^b_k(u)\|\le\delta_k.
\]
The threshold normalization explains this loss. For one component, let
$z$ be its local coefficient, $\xi$ its purchased readout, $\Gamma$ the
finite-volume column map, and $v$ the purchased label. The local pricing
and zero-scalar identities give $\xi=\vartheta_k^2z$ and $Kz=z$; the
purchased bias says
$\|\xi-\Gamma v\|_{\rm c}\le\vartheta_k^2\delta_k$, with
$\vartheta_k>0$. The recorded coefficient is
$\vartheta_k^{-2}K\Gamma v$. Write $\|\cdot\|_{\rm c}$ for the sum of
the absolute coordinate values. Since $K$ contracts the coordinate supremum
norm and that norm is bounded by $\|\cdot\|_{\rm c}$,
the bias gives
\[
 \|z-\vartheta_k^{-2}K\Gamma v\|_\infty
 \le\vartheta_k^{-2}\|\xi-\Gamma v\|_{\rm c}\le\delta_k.
\]
Taking the maximum over the finitely many components gives the displayed
array estimate. The weighted readout acts on each coordinate by the
threshold raised to its color-factor count. The simple-group low-grade
classification gives count two at every non-scalar coordinate; the scalar
coordinate of $z$ is its zero-jet evaluation, which equals the observable's
value at the identity and vanishes for these branches. Thus the weighted-readout definition gives
$\xi=\vartheta_k^2z$ coordinate by coordinate. The local identification
of this weighted readout, the zero-jet reading, and the purchased bias are
inputs here.
The selected clock tends to infinity, so $\delta_k\to0$: its logarithm is
minus one quarter of a quantity tending to infinity. Also, for each $\varepsilon>0$, eventually in $k$ and on an
ultrafilter-large set of $u$ there is $w=w(k,u)$ with
$\|w-w^b\|\le\varepsilon$ and $\mathsf R C^b_k(u)=T w$. The two
correction branches use their displayed coefficient triples exactly;
the principal coefficient is controlled only asymptotically. Finally,
at every sufficiently large fixed $k$ the ultrafilter limit of
$\mathsf R R_k(u)$ is the previously defined $W^b_k$.

Intersecting the two large sets at fixed $k$ gives
\[
 \|\mathsf R R_k(u)-T w^b\|
 \le \|\mathsf R\|\delta_k+\|T\|\varepsilon.
\]
The norm is continuous and the displayed closed ball is preserved by the
ultrafilter limit, so the same bound holds with $W^b_k$ on the left.
First let $k\to\infty$, then let $\varepsilon\downarrow0$: this proves
\eqref{eq:ex-branch-coefficient-input}. The purchased bias and the existence
of the finite-volume coefficient witnesses are retained inputs; this
finite-dimensional passage does not construct them.

The principal witness can be obtained from the local projection and the
canonical stress classification, subject to the following local inputs.
Assume the coefficient projection is linear on the finite sums below and
annihilates each forward difference of the localized trace and
antisymmetry currents. Assume also that the local principal coefficients
have zero scalar coordinate. The exact finite-volume antisymmetry identity
expresses each difference between transposed entries as a sum of those
forward differences. Their projected coefficients therefore form a
symmetric array. Its finite signed and real-structure average is a real
covariant symmetric tensor with zero scalar coordinate. The free stress
classification in Section~\ref{sec:ex-lattice-trace-free-rank} expresses it
as $T w$ for a triple $w$, giving the required witness. The trace-current
identity and the same vanishing give $\operatorname{tr}Tw=0$; evaluation on
a nonzero $01$ curvature jet yields $w_1+4w_2=0$. This exact local constraint
does not give a rate for the recorded field trace. The difference-projection
identity follows from the local affine-test identity: translating an affine test changes its potential by an infinitesimal
gauge term. The current's first variation at the identity vanishes: global
gauge invariance makes its restriction to each link an invariant linear
functional on the simple color algebra. Such a functional annihilates all
brackets, which span the nonabelian simple algebra, so it is zero. Differentiating gauge
invariance at the identity then shows that the Hessian annihilates an
infinitesimal gauge direction in either argument: the mixed derivative of
the gauge action is multiplied by the already vanishing first variation.
This also uses the local
realization of the gauge term on the nonwrapping support.
The Hessian therefore agrees on the original and translated affine tests.
The forward difference consequently has zero Hessian on each such test
direction paired with itself. Its value at the identity is also zero, and the
grade-four jet-reading identity makes every jet evaluation of its local
coefficient zero; the coefficient itself is zero. This application uses the affine translation
identity, local regularity and nonwrapping of its designated neighborhood,
Taylor order at least four, side length at least $40$, and the selected local
certificates for the antisymmetry currents and their differences. In the selected
construction the same finite-volume certificate supplies these local current and
difference bounds for every axis on a common ultrafilter-large set at each
sufficiently large scale. The existence of those choices and their local
bounds remains an input here; the finite-dimensional classification does
not itself produce them.

For the principal branch, the closeness of these witnesses has a useful
sequential reduction. Let $t_n\to\infty$ be any sequence of inverse
couplings and let the finite-volume data be admissible at each $t_n$, with
side length tending to infinity, the prescribed nonwrapping neighborhood,
and Taylor order at least four. The local comparison supplies two Hessian
estimates for each fixed pair of the finitely supported elementary test
directions of total degree at most four. Fixing a pair of tensor indices, let $Q_n$ be the corresponding
unmodified local observable and $\widetilde Q_n$ its principal
modification; their second derivatives on those directions differ by at most $16t_n^{-30}$, while the derivative of $Q_n$ differs
from the prescribed limiting bilinear form by at most $t_n^{-100}$. For
the first bound, write the Hessian of $Q_n-\widetilde Q_n$ in lattice
coordinates. If $r_i$ is the distance weight of coordinate $i$, each
elementary direction has $|v_i|\le4(1+r_i)^2$. Thus the product of
coordinates from the two directions is at most $16(1+\max(r_i,r_j))^4$.
The weighted Hessian coefficient sum controlling the modification dominates this polynomial
weight and is at most $t_n^{-30}$; summing the coordinate expansion gives
the factor $16$. The weighted coefficient estimate and the elementary
direction bound are the local inputs used here. The
canonical quadratic stencil has derivatives converging to that same form
as the side length grows. Thus the corresponding derivative of
$\widetilde Q_n$ and that of the canonical stencil have the same limit. These estimates use the local regularity and
comparison inputs and are not consequences of the later field limit.

The resulting Taylor-coefficient convergence follows as follows. A
first-order jet produces a lattice test direction which is a finite linear
combination of those elementary directions. Bilinearity and the preceding
bounds therefore give convergence of the second derivative on this test
direction paired with itself. The grade-four jet-reading identity expresses
the label evaluated at a jet as the observable's value at the identity
plus one half of this second derivative. For the principal observable the
identity value is zero, and the same identity for the canonical stencil
identifies the limit with the evaluation of the canonical stress label.
This uses the jet-reading and canonical-stencil identities under the
simple-group, regularity, nonwrapping and order hypotheses; they are
external inputs here. The jet evaluations separate the finite-dimensional
label space, so finitely many of them already separate it. Their joint
linear map has a continuous left inverse onto its image. Consequently the
unsymmetrized coefficient array $C_n$ tends entrywise to $T w^0$.
The limiting material array is fixed by the finite symmetry average, so
$\mathsf R C_n\to T w^0$. If $w_n$ is any associated coefficient witness
with $\mathsf R C_n=T w_n$, then $w_n\to w^0$. Indeed, the material coefficient map is injective by the curvature-jet
calculation above; choose a fixed linear left inverse on its finite-dimensional
range and apply it to the convergence.

This sequential statement yields the eventual uniform estimate required
above. Otherwise, for some $\varepsilon>0$, choose counterexamples with
$t_n\ge n$ and $\|w_n-w^0\|\ge\varepsilon$. Their admissible side length
$N_n$ satisfies $N_n>20(H(t_n)+10)$, with $H$ the neighborhood radius
defined below. Since $H(t_n)\to\infty$, the chosen data meet the sequential
hypotheses. The preceding convergence is a contradiction.
The local projection certificates, weighted Hessian bound, unmodified
Hessian comparison, jet-reading identities and neighborhood-size growth
are retained inputs. For the two correction branches the coefficient
identities follow directly from their finite-volume curvature products. The curvature
vanishes at the identity configuration, so for a lattice test direction
$v$,
\[
 \tfrac12D^2\langle F_{ij},F_{kl}\rangle_{1}[v,v]
 =\langle DF_{ij}(1)[v],DF_{kl}(1)[v]\rangle.
\]
On the affine test associated with a degree-zero radial jet, the two
linearized lattice curvatures are exactly the corresponding components of
that jet, provided the test stencil is nonwrapping. The grade-four
jet-reading identity then identifies this Hessian with $P_{ij,kl}$.
Taking the displayed finite sums gives $w^-=(0,0,1)$ and
$w^+=(0,\tfrac14,0)$ entry by entry, including their zero entries. Their
material tensors are invariant under the finite tensor-symmetry average,
so the same identities hold after that average. This uses the local
curvature-stencil and jet-reading identities as inputs and does not infer
an exact finite-volume equality from a continuum limit.

The same local identities give a useful, but slower, trace estimate. For
$b\in\{0,-\}$ the trace of $\mathsf R C^b_k(u)$ is zero on the relevant
large set: for $b=0$ this is the exact local constraint above, and for
$b=-$ it follows from the zero diagonal. Applying the finite trace map to
$\mathsf R(R_k(u)-C^b_k(u))$ and then taking the fixed-scale ultrafilter
limit gives
\[
 \|\operatorname{tr}W^b_k\|\le4\|\mathsf R\|\delta_k
\qquad(b=0,-).
\]
The factor four counts the diagonal entries. This does not prove the weighted trace limit:
for each fixed $p>0$, $\widehat\beta_k^p\delta_k\to\infty$. Indeed, for
$y=\log\widehat\beta_k$, the bound $\log y\le4y^{1/4}$ makes the logarithm
of this product at least $py-4\sqrt y\to\infty$. On the other hand,
$(\log\widehat\beta_k)^m\delta_k\to0$ for every fixed $m\in\mathbb N$, since its
logarithm is $m\log y-(\log y)^2/4$. The fast recording argument below
supplies a stronger estimate.

The moment-map convergence is a separate grade-four chart
input: the normalized transported lift tends to the core projection, the
scalar complement is annihilated by centering, and the chart expectations
converge. On the finite-dimensional label space, these statements give
operator-norm convergence. We spell out the selected two-scale passage
behind the normalized lift. Let $m_i$ be the selected initial level,
so $m_i\to\infty$, and let $k_i$ be its prescribed summit. The simple
group's grade-four reference projection is an isomorphism. To see this,
use the low-grade decomposition in Section~\ref{sec:ex-lattice-trace-free-rank}:
after the invariant linear terms vanish, a grade-four label is a constant
plus a quadratic expression depending only on its degree-zero radial jet.
Every such jet is realized by a constant-curvature affine Maxwell solution:
for prescribed antisymmetric $F_{ab}$ take
$A_b(x)=\frac12\sum_a F_{ab}x_a$. Its curvature is $F$, and all curvature
derivatives, including the Maxwell divergence, vanish. A label in the free
equation ideal therefore vanishes on every degree-zero jet and hence on every jet. The
standard table identifies the kernel of its projection with that equation
ideal, an earlier free-table input. The kernel is zero, and the projection
is onto by construction. In particular $\mathcal O\ne0$: on an affine
Maxwell solution with $F_{01}=v$, $F_{10}=-v$ and all other curvature
components zero, its unprojected label evaluates to $\tfrac12\langle v,v\rangle>0$
for any nonzero color vector $v$. Injectivity preserves its nonvanishing.
All three material tensors are nonzero: every $w^b\ne0$, and both the
material coefficient map and $\pi_0$ are injective. In particular the
false material tensor is nonzero although its trace vanishes by
\eqref{eq:ex-branch-three-traces}.
Write $\ell=\pi_0^{-1}$. Let
$Y_i:\mathcal F_4^0\to E_4$ be the selected label
lift, and set $\varpi_i=\beta_Z(m_i)r_i^4>0$. For each fixed
$z\in\mathcal F_4^0$, the incoming local chart estimates and the scale
comparison below give
\begin{align}
 \varpi_i^{-1}B_{m_i}Y_i z&\longrightarrow K\ell z,
 \label{eq:ex-branch-lift-initial}\\
 \|B_{m_i}Y_i z-B_{k_i}Y_i z\|
 &\le \varepsilon_i\|B_{k_i}Y_i z\|,\quad
 \varepsilon_i\longrightarrow0,
 \label{eq:ex-branch-lift-stability}\\
 \frac{\beta_Z(m_i)}{\beta_Z(k_i)}&\longrightarrow1.
 \label{eq:ex-branch-lift-beta-ratio}
\end{align}
The first convergence comes from the fixed-grade local lift estimate after
division by $\varpi_i$: for each fixed $z$ its error is bounded by a constant
depending on $z$ times
$\exp[-\tfrac12(\log\log\widehat\beta_{m_i})^2]$, which tends to zero.
We spell out its finite-volume normalization. At a fixed initial level $m$,
write $\mathcal N$ for the column map normalized by $a_m^{-4}$,
$\mathcal T$ for the local chart map, and let $\alpha$ denote the displayed exponential error. The
selected chart inputs give
$\mathcal N u_{\rm ref}=\mathcal T u_{\rm ref}=e_0$ and, for every $y$,
\[
 \|\mathcal Ny-\mathcal Ty\|_{\rm c}\le\alpha\|\mathcal Ny\|_{\rm c},\qquad
 \|\mathcal Ny-\mathcal Ty\|_{\rm c}\le\tfrac12\|\mathcal Ny\|_{\rm c}.
\]
The second estimate and the triangle inequality give
$\|\mathcal Ny\|_{\rm c}\le2\|\mathcal Ty\|_{\rm c}$.
For any $v$, take $y=v-(\mathcal Tv)_0u_{\rm ref}$. Then
$\mathcal Ty=K\mathcal Tv$ and $K\mathcal Ny=K\mathcal Nv$, so the
coefficient norm's contraction under $K$ gives
\[
 \|K(\mathcal Nv-\mathcal Tv)\|_{\rm c}\le2\alpha\|K\mathcal Tv\|_{\rm c}.
\]
At grade four the non-scalar coordinates of every label have the
required homogeneous grade, by the low-grade simple-group classification.
For the actual finite-volume lift $L$ of a reference label $v$, the
selected lift identity is
$K\mathcal TLv=\beta_{m,u}\vartheta_m^2 h^4Kv$, where
$\beta_{m,u}>0$, $\vartheta_m>0$, $h\ge0$ and the physical radius is
$a_mh$. The finite-volume counterpart of $B_mLv$ is
$\vartheta_m^{-2}a_m^4K\mathcal NLv$. Substituting $Lv$ into the preceding
inequality yields
\[
 \bigl\|\vartheta_m^{-2}a_m^4K\mathcal NLv-
       \beta_{m,u}(a_mh)^4Kv\bigr\|_{\rm c}
 \le2\alpha\,\beta_{m,u}(a_mh)^4\|Kv\|_{\rm c}.
\]
At fixed $i$ the selected lifts and couplings have their common
separation-ultrafilter limits, and the selected physical radius is $r_i$.
The steep transport property gives the limit of the finite-volume
transport applied to those moving lifts. Passing to these limits gives the
stated initial-level estimate; no separate limit of $\mathcal T$ is needed.
The two finite-volume chart inequalities, lift identity and selected
limits are the inputs used in this reduction.

The second estimate is the comparison over the actual interval from the
initial level to its summit. Its coefficient-size error is
\[
 \operatorname{part}(m_i)=
 \exp\!\left[-\tfrac25(\log\log\widehat\beta_{m_i})^2\right].
\]
The norm is any fixed norm on $E_4$; equivalence with the coefficient norm
transfers the estimate with $\varepsilon_i$ equal to a fixed
positive multiple of $\operatorname{part}(m_i)$. This tends to zero because
the hatted clock at the cofinal initial levels tends to infinity. The
finite-volume inputs behind these comparisons belong to the selected
history; they do not follow from a generic chart.

The interval comparison follows from the actual one-step column bound and
the scheduled sum budget. At one step, let $w$ and $v$ be the finite-volume
column values at the earlier and later levels, with positive thresholds
$a,b$. The selected one-step input gives $b/a\le2$ and a comparison value
$X$ with
\[
 \|w-X\|_{\rm c}\le g\|Kv\|_{\rm c},\qquad
 KX=(a/b)^2Kv.
\]
Here $g$ is the step's nonnegative column error; the second identity uses
the grade-four homogeneity just noted. Thus
\[
 \|a^{-2}Kw-b^{-2}Kv\|_{\rm c}
 \le a^{-2}g\|Kv\|_{\rm c}
 \le 4g\|b^{-2}Kv\|_{\rm c}.
\]
This is the one-step estimate with coefficient $4g$; the threshold and
column comparison are selected-construction inputs. To see the
accumulation, fix a finite-volume label and write $u_j$ for its transported
value at level $j$ between an initial level $m$ and its summit $k$. Put $N=k-m>0$ and
$x=\operatorname{part}(m)\le1$. The selected one-step and schedule inputs
state that
\[
 \|u_j-u_{j+1}\|_{\rm c}
 \le \frac{x}{2N}\|u_{j+1}\|_{\rm c},\qquad m\le j<k.
\]
Backward induction gives, for every $m\le j\le k$,
\[
 \|u_j-u_k\|_{\rm c}
 \le \frac{x(k-j)}N\|u_k\|_{\rm c}.
\]
Indeed, if this is known at $j+1$, the triangle inequality and $x\le1$
give $\|u_{j+1}\|_{\rm c}\le2\|u_k\|_{\rm c}$. Add the one-step error
to the induction bound to obtain the next case; the case $j=k$ is zero.
For fixed $m,k$ and label, the finite-volume estimates hold together for
the finitely many levels on an ultrafilter-large set. Passing to the
fixed-scale transport limits gives the stated comparison. The threshold and column comparison and the schedule hypotheses remain
selected-construction inputs; the induction does not establish them.

For the scale ratio, the selected finite-volume running estimate says
that for every $\epsilon>0$, eventually in $j$,
\[
 |\beta_{j,u}-\widehat\beta_j|\le\epsilon s_j
\]
on an ultrafilter-large set of finite-volume parameters $u$. At each fixed
$j$, the same ultrafilter sends $\beta_{j,u}$ to $\beta_Z(j)$; passing to
the closed interval gives the same estimate for $\beta_Z(j)$. The affine
clock satisfies $\widehat\beta_j\ge b s_j$ for its positive slope $b$.
For any $\delta>0$, choose $\epsilon=b\delta/2$ and divide by the positive
clock. This proves $\beta_Z(j)/\widehat\beta_j\to1$. The running estimate
and the fixed-scale limit are upstream inputs. The summit clock estimate follows from the first-crossing and
block-size properties of the selected schedule. Write $g_j$ for the gap
from level $j$ to its summit and $b$ for the fixed slope constant. Set
$\operatorname{Gate}(x)=(\log x)^{48}$. For $x\ge1$, the inequality
$t\le e^t$ with $t=(\log x)/96$ gives
$(\log x)^{48}/x\le96^{48}x^{-1/2}$, so the gate ratio tends to zero.
The selected clock tends to infinity.

The summit is the first later level at which the scale gap reaches the
gate. Let $q$ be its preceding level and put $x=\widehat\beta_j$. By the
first-crossing property, $0\le s_q-s_j\le\operatorname{Gate}(x)$.
The affine clock identity then gives
$\widehat\beta_q\le x+b\operatorname{Gate}(x)\le2x$ eventually. The
selected schedule has $L_q\ge1$ and bounds $L_q$ by
$\log\widehat\beta_q$; its scale increment is
$s_{q+1}-s_q=\log L_q$. Since $\log L_q\le L_q$, for large $x$,
\[
 s_{q+1}-s_q\le\log\widehat\beta_q
 \le\log(2x)\le\operatorname{Gate}(x).
\]
Adding the pre-crossing gap and this last increment gives
$0\le g_j\le2\operatorname{Gate}(\widehat\beta_j)$. Consequently,
\[
 \widehat\beta_{\mathrm{summit}(j)}
   =\widehat\beta_j+b g_j,\qquad
 0\le\frac{g_j}{\widehat\beta_j}
 \le 2\frac{\operatorname{Gate}(\widehat\beta_j)}{\widehat\beta_j}
 \longrightarrow0.
\]
The first-crossing definition and block-size schedule bound are the inputs
used here. Dividing the identity by the positive $\widehat\beta_j$ yields
$\widehat\beta_{\mathrm{summit}(j)}/\widehat\beta_j\to1$, and hence the
same limit along $j=m_i$. Since both levels are cofinal and the
denominators are positive,
\[
 \frac{\beta_Z(m_i)}{\beta_Z(k_i)}=
 \frac{\beta_Z(m_i)/\widehat\beta_{m_i}}
 {\bigl(\beta_Z(k_i)/\widehat\beta_{k_i}\bigr)
  \bigl(\widehat\beta_{k_i}/\widehat\beta_{m_i}\bigr)}
 \longrightarrow1.
\]
Thus \eqref{eq:ex-branch-lift-beta-ratio} follows from these two explicit
scale estimates.

The same bounds also explain the scheduled budget used in the backward
induction. Put $x=\widehat\beta_m$ and $N=k-m$, where $k$ is the summit.
Every scale increment is at least one, so the gap bound gives
$N\le2(\log x)^{48}$. For all large $m$, the summit clock is at most
$2x$, and the scheduled column factor $p_j$ at each step satisfies
$p_j\le2\log x$. Let
$c_j=\exp[\tfrac12(\log\log\widehat\beta_{j+1})^2]$; clock monotonicity
gives $c_j\ge\exp[\tfrac12(\log\log x)^2]$ at these levels. The
one-step coefficient in the selected construction is
$8(1+p_j)/c_j$. To bound it by $\operatorname{part}(m)/(2N)$, it is
enough that
\[
 16N(1+2\log x)\le\exp[\tfrac1{10}(\log\log x)^2].
\]
Using the bound on $N$, the left side is at most
$32(1+2\log x)(\log x)^{48}$. With $y=\log\log x$, this is at most
$96e^{49y}$ for large $x$, whereas the ratio of the right side to
$96e^{49y}$ tends to infinity. This proves the scheduled budget from the
clock, block and column bounds. The underlying threshold and column
comparison remains an input; it is not a consequence of this numerical
calculation.

To combine them, put $X_i=\varpi_i^{-1}B_{m_i}Y_i z$ and
$Z_i=\varpi_i^{-1}B_{k_i}Y_i z$. The stability estimate gives
$\|X_i-Z_i\|\le\varepsilon_i\|Z_i\|$. Since $X_i$ converges, it is
bounded; eventually $\varepsilon_i\le1/2$, so
$\|Z_i\|\le2\|X_i\|$. Consequently $Z_i-X_i\to0$. Multiplication by
\eqref{eq:ex-branch-lift-beta-ratio} changes the denominator from
$\beta_Z(m_i)r_i^4$ to $\beta_Z(k_i)r_i^4=\omega_i$, giving
\begin{equation}
 \omega_i^{-1}B_{k_i}Y_i z\longrightarrow K\ell z.
 \label{eq:ex-branch-lift-summit}
\end{equation}
The same $Y_i$, $r_i$, transport and history occur at both levels; the ratio is
not a comparison between independently selected charts.

Now put
\[
 D_i=\omega_i^{-1}B_{k_i},\qquad
 \mathcal A_i=D_iY_i\pi_0+(1-K):E_4\longrightarrow E_4.
\]
The limit \eqref{eq:ex-branch-lift-summit}, $\ell\pi_0=1$, and a finite
basis give $\mathcal A_i\to1$ in operator norm; here $K$ is the core projection on
the fixed label space. Also $\mathsf M_i$
annihilates the scalar complement $1-K$, and eventually
$\pi Y_i=U_i$. Thus, on $x\in E_4$, the chart identity gives
\[
 \mathsf M_i \mathcal A_i x
 =\pi_0x-\langle U_i\pi_0x\rangle\mathbf1_4^0
 \longrightarrow \pi_0x-\langle\pi_0x\rangle\mathbf1_4^0
 =\mathsf M_\infty x.
\]
The expectation convergence is the one-point case of the selected chart
correlation limit. For a fixed standard field $z$, its one-point function
at the origin is $\langle z\rangle_0$; the rescaled interacting one-point
function at the origin is $\langle U_i z\rangle$. The selected chart
convergence at this one-point configuration therefore gives
$\langle U_i z\rangle\to\langle z\rangle_0$. It uses the chart's
one-point convergence and the two table hierarchies as inputs. Eventual
transport of the unit is another selected-chart input. A finite basis
upgrades the displayed pointwise convergence to operator-norm convergence.
Eventually $\|\mathcal A_i-1\|<1/2$, so the geometric series gives
$\|\mathcal A_i^{-1}\|\le2$ and
$\|\mathcal A_i^{-1}-1\|\le2\|\mathcal A_i-1\|\to0$.
Since $\mathsf M_i=(\mathsf M_i\mathcal A_i)\mathcal A_i^{-1}$,
this proves \eqref{eq:ex-branch-moment-input}. The normalized local lift
and one-point chart convergence remain upstream dependencies. Neither limit alone
supplies a trace rate.

\begin{proposition}[Three branch field limits]
\label{prop:ex-branch-moments}
Assume \eqref{eq:ex-branch-coefficient-input}--
\eqref{eq:ex-branch-moment-input}, the chart identity and projection
equivariance above, and eventual equivariance of $B_k$ under every signed
motion and the real-structure involution. Then
\[
 \mathsf P_i t^0_{k_i}\longrightarrow M(w^0),\qquad
 \mathsf P_i t^-_{k_i}\longrightarrow M(w^-),\qquad
 \mathsf P_i t^+_{k_i}\longrightarrow M(w^+)
\]
in the finite product $(\mathcal F_4^0)^{I\times I}$.
\end{proposition}
\begin{proof}
The finite symmetry average commutes with symmetrization. For sufficiently
large $k$, the selected transport intertwines every signed motion and the
real-structure involution: this follows by taking its defining fixed-scale
separation limit of the corresponding finite-volume transport identities.
The core projection also commutes: the scalar label is fixed and its
coordinate is preserved, up to conjugation under the real structure. Using
equivariance and the chart identity gives, eventually,
\begin{equation}
 (\mathsf P_i t^b_{k_i})_{ac}=
 \tfrac12\bigl(\mathsf M_i W^b_{k_i}(a,c)
                 +\mathsf M_i W^b_{k_i}(c,a)\bigr).
 \label{eq:ex-branch-chart-array}
\end{equation}
Indeed, apply \eqref{eq:ex-branch-chart-identity} to each term of the
finite symmetry average; the scalar signs and coefficients are real, so
commute with both maps. Since $k_i\to\infty$,
$W^b_{k_i}(a,c)\to W^b(a,c)$. For any convergent $v_i\to v$ in $E_4$,
\[
 \|\mathsf M_i v_i-\mathsf M_\infty v\|
 \leq \|\mathsf M_i-\mathsf M_\infty\|\,\|v_i\|
       +\|\mathsf M_\infty\|\,\|v_i-v\|\longrightarrow0.
\]
Apply this to both entries in \eqref{eq:ex-branch-chart-array}. The limit
is $\operatorname{sym}(\mathfrak c_0\pi_0\mathcal S(w^b))$, which is
$M(w^b)$ by symmetry and zero expectation. There are only sixteen entries,
so coordinate convergence is convergence of the entire tensor. This proves
all three limits without exchanging an infinite sum and a limit.
\end{proof}

\subsection{The fast trace from localized recordings}
\label{subsec:ex-branch-fast-trace}

We now derive the raw fast-trace statement used in the weighted limit. Let
$s_k\to\infty$ be the logarithmic scale of the selected schedule and
$a_k=e^{-s_k}$ for the mesh above. A sequence is \emph{fast} if its norm is
$O(e^{-\eta s_k})$ for some $\eta>0$; a scalar sequence or operator family
is \emph{subexponential} if its norm is $O(e^{\epsilon s_k})$ for every
$\epsilon>0$. A finite sum of fast sequences is fast, with the minimum of
finitely many positive decay exponents.

We make the local recording input explicit. Write $E_3$ and $\mathcal F_3^{\mathrm{int}}$
for the finite-dimensional grade-three label and field spaces of the same
interacting table, and $\pi_d:E_d\to\mathcal F_d^{\mathrm{int}}$ for its grade-$d$
projection, with $\pi_4=\pi$.
For each admissible separated test $\sigma$, the table gives linear functionals
$L^{\sigma}_d(z):\mathcal F_d^{\mathrm{int}}\to\mathbb C$ for $d=3,4$ and $z$ near the origin.
They are continuous in $z$ and $L^{\sigma}_d(0)=\lambda^{\sigma}_d$. The functionals
$\lambda^{\sigma}_4$ separate $\mathcal F_4^{\mathrm{int}}$; hence finitely many already do so.
This follows from the field-separation and smoothness properties of the
selected hierarchy: approximate a separated configuration by generic grid
configurations at increasingly fine scales. Subtracting the first grid
address places the insertion at the origin without changing relative
positions. Translation invariance preserves the correlation, and spectator
fields have representative tokens, so these are allowed tests. By continuity,
a field annihilated by them has all separated correlations zero and is zero
by field separation. These hierarchy and
grid-density properties are retained inputs; finite dimensionality alone
does not imply separation.

Every grade-three field is a multiple of the unit. Indeed, each color-valued
jet-coordinate factor has grade at least two, so a grade-three label has at
most one factor. Its invariant linear part vanishes by the low-grade color
argument of Section~\ref{sec:ex-lattice-trace-free-rank}; only the scalar
label remains. The field projection is onto and sends the scalar label to
a nonzero reference-scale multiple of the unit, so the same is true of
$\mathcal F_3^{\mathrm{int}}$. Finally, inserting the unit leaves a separated
correlation independent of its insertion point. Thus $L^{\sigma}_3(z)$ is
independent of $z$ in a sufficiently small separated neighborhood. The
grading, quotient projection, unit and field-separation statements here
are the ones belonging to the same selected table; they are not inferred
from the branch limit.

The incoming recording estimates are the following. For every test $\sigma$
there is a radius $c_\sigma>0$. For each $d\in I$ the localized current
$Q_d^{[r]}$ and its forward difference have recorded fields
$v_{d,k}\in\mathcal F_3^{\mathrm{int}}$ and $w_{d,k}\in\mathcal F_4^{\mathrm{int}}$. For every branch
$b\in\{0,-\}$, the diagonal tensor recordings are the already defined
$\pi x^b_k(a,a)\in\mathcal F_4^{\mathrm{int}}$. Their finite-volume tested correlations,
denoted respectively by $F^{\sigma}_{d,k,u}(j)$, $G^{\sigma}_{d,k,u}(j)$ and
$T^{\sigma,b}_{a,k,u}$, satisfy, with fast nonnegative errors, the following
estimates. For all sufficiently large $k$ and each integer $j$ with
$|j|a_k\le c_\sigma$, the appropriate inequalities hold for $u$ in an
ultrafilter-large set. A common radius is obtained by taking the minimum
of the finitely many radii for these current, difference and tensor
insertions; the corresponding eventual scale cutoff is the maximum of
finitely many cutoffs:
\begin{align}
 |F^{\sigma}_{d,k,u}(j)-L^{\sigma}_3(ja_ke_d)v_{d,k}|&\le E^{\sigma}_{d,k},
 \label{eq:ex-branch-book-current}\\
 |G^{\sigma}_{d,k,u}(j)-L^{\sigma}_4(ja_ke_d)w_{d,k}|&\le D^{\sigma}_{d,k},
 \label{eq:ex-branch-book-difference}\\
 |T^{\sigma,b}_{a,k,u}-\lambda^{\sigma}_4(\pi x^b_k(a,a))|&\le H^{\sigma,b}_{a,k}.
 \label{eq:ex-branch-book-tensor}
\end{align}
Here $e_d$ is the $d$th coordinate vector; $E,D,H$ are fast sequences,
possibly with test- and component-dependent exponents. The same separation
ultrafilter is used in every line. The last line is at zero displacement.
The finite-volume correlation is linear in its inserted observable, and
for each $\sigma$ the forward-difference identity holds for all integers
$j$ on an ultrafilter-large set, for all sufficiently large $k$:
\begin{equation}
 G^{\sigma}_{d,k,u}(j)=a_k^{-1}
 \bigl(F^{\sigma}_{d,k,u}(j+1)-F^{\sigma}_{d,k,u}(j)\bigr).
 \label{eq:ex-branch-recorded-difference}
\end{equation}
Here $a_k^{-1}=a_k^{-4}/a_k^{-3}$, the ratio of the two grade
normalizations; translation covariance identifies the displaced current values.
The localized current and its difference are recorded using the same finite-volume choices
and volumes. The construction supplies these estimates using its local
recording, selected trace-current localization, field hierarchy and
point-limit inputs. The estimates, low-grade unit statement and separating
property use the upstream table and hierarchy inputs just stated; the
finite-volume price production behind the estimates is not proved in this
subsection. They are more primitive than the assertion that either branch
trace is fast.

The following root-price mechanism identifies a more primitive sufficient
input for the three recording estimates. Fix one of the finitely many
insertions above, its grade $d\in\{3,4\}$, and a separated test $\sigma$.
At scale $k$ and finite-volume parameter $u$, let $A_{k,u}$ be the actual
inserted observable and $O_{k,u}(v_{k,u})$ the observable made from its
selected local label $v_{k,u}\in E_d$. The actual insertion already includes
the grade-$d$ factor $a_k^{-d}$ and transport to the comparison volume;
$O_{k,u}$ is the corresponding reference-scale observable. The
following finite-dimensional selection explains how the same limiting label is obtained for all tests.
The purchased finite-volume recording supplies a coefficient
$\xi_{k,u}$, a label $v_{k,u}$, and the actual column map $C_{k,u}$ on an
ultrafilter-large set, with
\begin{equation}
 \|\xi_{k,u}\|_{\mathrm c}\le1,\qquad
 \|\xi_{k,u}-C_{k,u}v_{k,u}\|_{\mathrm c}\le1,
 \qquad
 \|v_{k,u}\|\le M_{d,k}\|C_{k,u}v_{k,u}\|_{\mathrm c}.
 \label{eq:ex-branch-recording-compactness}
\end{equation}
Here $\|\cdot\|_{\mathrm c}$ is the fixed coefficient norm, and
$M_{d,k}=\exp(\ell_k+ds_k)$ and
$\ell_k=\sum_{j<k}\log\zeta_{j+1}$ for the calibrated column factors
$\zeta_{j+1}\ge1$. The previously established column estimate gives
$\ell_k=o(s_k)$, so $e^{\ell_k}$ is subexponential. The middle bound
is a convenient weakening of the purchased bias bound
$\vartheta_k^2\exp[-\tfrac14(\log\log\widehat\beta_k)^2]$, using
$\vartheta_k\le1$. It is not a replacement for the root-price purchase
below. The inverse bound follows by multiplying the calibrated one-step factors
$\zeta_{j+1}L_j^d$ from the reference level, where the column is the identity.
All factors are at least one, $\sum_{j<k}\log L_j=s_k$, and the label norm
is bounded by the coefficient norm. The calibration, reference identity
and common choice of column map and label remain inputs.

When constructing a history from these purchased data, for each fixed
sufficiently large $k$ choose $v_{k,u}$ once together with its full
purchase certificate on the ultrafilter-large set where that certificate
is available, and define it arbitrarily elsewhere. In particular,
\eqref{eq:ex-branch-recording-compactness} holds on that set. By the
triangle inequality its norm is at most
$2M_{d,k}$ on that set. The closed ball of that radius is compact because
$E_d$ is finite dimensional, so the ultrafilter has a limit $x_k\in E_d$
in the ball. The choice is made before any test is selected; this one
$x_k$ works for every fixed test and satisfies $\|x_k\|\le2M_{d,k}$.
Since $\ell_k=o(s_k)$, for every $\epsilon>0$ this also gives
$\|x_k\|=O(e^{(d+\epsilon)s_k})$; the bound is grade-dependent and need
not be subexponential. For an already selected history, retain its original labels and their
existing ultrafilter limit instead: the compactness construction by itself
does not identify a newly chosen limit with that history's limit. The
root-price step below uses the labels and limit from one and the same
history.

For the root-price step assume:
\begin{enumerate}
\item for this test, some $c_\sigma>0$ and $B_\sigma\ge0$ satisfy, for
 every sufficiently large $k$, axis $a\in I$ and integer $\ell$ with
 $|\ell|a_k\le c_\sigma$, on an ultrafilter-large set of $u$,
\[
 |\mathcal C^{\sigma}_{k,u}(\ell,a;A_{k,u})-
   \mathcal C^{\sigma}_{k,u}(\ell,a;O_{k,u}(v_{k,u}))|
 \le B_\sigma E_k,
\]
 where $E_k$ is fast and $\mathcal C$ is the ordinary finite-volume
 tested correlation, linear in the insertion;
\item for each sufficiently large fixed $k$, axis $a$, integer $\ell$ with
 $|\ell|a_k\le c_\sigma$, and fixed $v\in E_d$, the correlation of
 $O_{k,u}(v)$ converges along the same ultrafilter to
 $L^{\sigma}_d(\ell a_ke_a)\pi_d v$.
\end{enumerate}
The first assumption applies the rooted price to the actual-minus-local
insertion, using selected funding, interior and locality estimates. The
second is the fixed-label point limit in the same table. To apply it to the moving labels,
choose a basis $(e_j)$ of $E_d$ and write
$v_{k,u}=\sum_j c_j(u)e_j$, with $c_j(u)\to c_j$ along the ultrafilter.
Linearity of the local insertion and correlation gives a finite sum of
$c_j(u)$ times the correlations for the fixed basis vectors. The product
and finite-sum limit rules therefore give
\[
 \mathcal C^{\sigma}_{k,u}(\ell,a;O_{k,u}(v_{k,u}))
 \longrightarrow L^{\sigma}_d(\ell a_ke_a)\pi_d x_k.
\]
This uses the existing fixed-scale label limit, including for an already
chosen history. Fix a lower scale $h$ at least the test's native scale. A comparison scale is open when $h\le j$ and $b_j^2\le a_h/a_j$, where $b_j$
is its cover size. Let $\operatorname{basin}(h)$ be the first open scale and
choose $j=\max\{\operatorname{basin}(h),\operatorname{ref}(d)\}$. Such a scale exists: the cover-growth and mesh bounds give
$a_{j+1}b_{j+1}^2\le(2/3)a_jb_j^2$. Iterating gives
$a_jb_j^2\le(2/3)^{j-h}a_hb_h^2\to0$ for $j\ge h$. Opening persists at
finer scales, so this choice of $j$ is open. At each allowed later level $k>j$, the
finite-volume purchase gives the actual-minus-local insertion the rooted price
\begin{equation}
 P_{d,h,j}(k)=
 \frac{2\,\operatorname{stake}_{h,j}\,a_j^{-d}}
      {\varphi_{d,j}}\,
 \widehat\beta_k^4\,
 \exp\!\left[-\frac{s_k-s_j}{8}\right].
 \label{eq:ex-branch-native-root-price}
\end{equation}
Here $\operatorname{stake}_{h,j}=\exp[(a_h/a_j)^4\mathcal T_j+5]$,
where $\mathcal T_j$ is the selected per-site treasury, and
$\varphi_{d,j}=\widehat\beta_j^{-\nu_d}$, where $\nu_d$ is the fixed
reference-label degree bound plus four. Both are positive and fixed once
$d,h,j$ are chosen. This is the actual purchased price. Since the hatted clock is affine in $s_k$, its
fourth power is subexponential: for $s\ge0$ and $\epsilon>0$, the
exponential series gives $s^4\le24\epsilon^{-4}e^{\epsilon s}$. Thus, for
every $0<\eta<1/8$, $P_{d,h,j}(k)=O(e^{-\eta s_k})$ and is fast. The containing-window radius
has an explicit bound. Put $r(x)=\lceil(\log x)^2\rceil$,
$q(x)=r(x)^{60}+1100$, and
\[
 H(x)=1800(1001r(x))+4(1000r(x)+2),\qquad N_0=4000001^4.
\]
The selected radius is $R_k=R(\widehat\beta_k)$, where
\[
 R(x)=\max\{H(x)+18+q(x)(4r(x)+2)+2r(x),\,
              2((N_0+q(x))(4r(x)+2)+10r(x))+1\}.
\]
For $x\ge e$, $r(x)\le1+(\log x)^2$, so this expression is at most a
constant times $(1+\log x)^{122}$. Also $\log x\le32x^{1/32}$, whence
$R(x)=O(x^{122/32})=o(x^4)$. Consequently
$R_k\le\widehat\beta_k^4$ eventually and $a_kR_k\to0$. This radius
includes the support enlargement and is distinct from the principal
localization radius $\varrho_k$: in particular $R(x)\ge H(x)+18$, so the
local regularity neighborhood of radius $H(x)+1$ lies in the window.
The local symmetry, measurability and size bounds needed for the interior
certificate remain inputs. The selected geometric tail is
$\operatorname{tail}_k=30a_k(J_k+5)$, where $J_k\ge0$ is the schedule's
integer clock. The schedule gives $a_{k+1}\le a_k/3$ and
$J_{k+1}\le J_k+1$. Hence
\[
 \operatorname{tail}_{k+1}
 \le10a_k(J_k+6)\le\tfrac23\operatorname{tail}_k,
\]
so the tail tends to zero geometrically. These two limits eventually give
the geometric room inequality
\[
 2(a_kR_k+\operatorname{tail}_k)<a_h.
\]
It is this inequality that allows the localized observable to sit in the
actual fixed test chamber before applying the purchased root bound.
To connect that bound to the first assumption, put
$f_{k,u}=A_{k,u}-O_{k,u}(v_{k,u})$. The purchase and localization say that,
on an ultrafilter-large set, this difference is interior to the chamber at
level $h$ and its root-test size is at most $P_{d,h,j}(k)$. The separated
spectator estimate supplies a constant $B_\sigma\ge0$, independent of $k$,
and a possibly smaller displacement radius, such that
\[
 \bigl|\mathcal C^{\sigma}_{k,u}(\ell,a;f_{k,u})\bigr|
 \le B_\sigma\operatorname{root}_{h,u}(f_{k,u}).
\]
The spectator input places the $n$ translated factors in distinct cells
and bounds each root-test size by a common $A\ge0$, uniformly in $k$ and
displacements of size at most $a_h/8$. The rooted moment inequality gives
the bound with $B_\sigma=A^n$; its integrability and separation conditions
are part of that input. To see the room for the local representatives,
choose a finer level $j'\ge h+1$ with $24a_{j'}<a_h$ at which the
spectator input supplies a common bound for its finitely many tokens and
frames. It follows from the centered-representative hypothesis, using
exponent and error allowances $1$, then summing the finitely many bounds
and intersecting their eventual sets. The displacement and three fine-mesh
widths together are less than $a_h/4$, inside the coarse cell's one-third
margin. For sufficiently large period
the spectator addresses and the insertion address occupy distinct cells.
By correlation linearity, the first assumption follows with
$E_k=P_{d,h,j}(k)$ after intersecting the relevant large sets. The local
representative and its root bound still use the spectator input. The cover-growth bound, purchase, localization, spectator estimate and the
displayed mesh and integer-clock bounds remain inputs from the selected
construction; this argument does not establish those purchases.

For each fixed $k$ and displacement, the moving-label limit makes the
local correlation differ from its limit by at most $a_k$ on an ultrafilter-large
set. Intersect that set with the root-price set and use the triangle
inequality. The resulting recording error is at most
$B_\sigma E_k+a_k$, which is fast because $a_k=e^{-s_k}$. This proves
the corresponding instance of
\eqref{eq:ex-branch-book-current}--\eqref{eq:ex-branch-book-tensor}
from the displayed root-price and point-limit inputs. It does not construct
the underlying purchased root bound. The radius can be reduced so that all moved
test configurations remain separated.

\begin{proposition}[Localized current and raw trace decay]
\label{prop:ex-branch-raw-fast}
Assume the recording and test properties above, and the exact
finite-volume branch formulas. Then for each $d$ the recorded difference
$w_{d,k}$ is fast, and so are both unaveraged projected traces
\[
 y^0_k=\sum_a\pi x^0_k(a,a),\qquad
 y^-_k=\sum_a\pi x^-_k(a,a).
\]
\end{proposition}
\begin{proof}
We first record the finite-dimensional detection step. Choose finitely many
$\lambda^{\sigma}_4$ whose joint map is injective. If for each such $\sigma$ linear
functionals $\Lambda^{\sigma}_k$ converge pointwise to $\lambda^{\sigma}_4$, finite
dimensionality gives convergence in operator norm. If $\mathcal F_4^{\mathrm{int}}=\{0\}$ the conclusion is immediate. Otherwise the joint
limiting map has a positive minimum on the unit sphere. Consequently, for all large
$k$ and every $z\in\mathcal F_4^{\mathrm{int}}$,
\begin{equation}
 \|z\|\le C\max_\sigma|\Lambda^{\sigma}_k z|.
 \label{eq:ex-branch-moving-detector}
\end{equation}
Thus fast bounds on these finitely many moving tests imply a fast field
bound, even when their exponents depend on $\sigma$.

Fix $d$ and a test $\sigma$. Decrease the decay exponent if necessary so that
$0<\eta<1$ and $E^{\sigma}_{d,k}=O(e^{-\eta s_k})$. Choose
$0<\tau<\eta$ and set
\[
 n_k=\left\lceil e^{(1-\tau)s_k}\right\rceil,\qquad
 h_k=a_kn_k.
\]
The ceiling bound gives
\[
 e^{-\tau s_k}\le h_k
 \le e^{-\tau s_k}+e^{-s_k}.
\]
Thus $n_k\ge1$, $h_k\to0$, and
$E^{\sigma}_{d,k}/h_k=O(e^{-(\eta-\tau)s_k})$. For large $k$ all the
points $ja_ke_d$, $0\le j\le n_k$, lie inside the recording
neighborhood. Define the moving functional
\[
 \Lambda^{\sigma}_k=\frac1{n_k}\sum_{j=0}^{n_k-1}L^{\sigma}_4(ja_ke_d).
\]
Continuity at zero implies $\Lambda^{\sigma}_k z\to\lambda^{\sigma}_4z$ for every
fixed $z$: every term in the average lies within distance $h_k$ of
zero. At each fixed $k$ intersect the finitely many ultrafilter-large sets
needed for the two endpoint current estimates, the $n_k$ difference
estimates, and the exact difference identity. This intersection is
nonempty. On it the grade-three unit property gives a common value
$C_{d,k}^{\sigma}=L^{\sigma}_3(0)v_{d,k}$ at both endpoints. Telescoping
\eqref{eq:ex-branch-recorded-difference} and using the recording errors
therefore yields
\[
 |\Lambda^{\sigma}_k w_{d,k}|
 \le D^{\sigma}_{d,k}+\frac{2E^{\sigma}_{d,k}}{a_kn_k}.
\]
The right side is fast. This construction may use a different $n_k$ for
each of the finitely many tests; applying
\eqref{eq:ex-branch-moving-detector} proves that $w_{d,k}$ is fast.
The finite intersection at fixed $k$ is essential: no common infinite
intersection of ultrafilter-large sets is used.

For the false branch the sum of the finite-volume diagonal insertions is
exactly zero. For a fixed separating test, intersect the four sets in
\eqref{eq:ex-branch-book-tensor} with that exact identity and use
linearity. The triangle inequality gives
$|\lambda^{\sigma}_4(y^-_k)|\le\sum_a H^{\sigma,-}_{a,k}$, which is fast. The
fixed-test version of \eqref{eq:ex-branch-moving-detector} proves the
false trace is fast.

For the principal branch, its finite-volume trace is exactly the sum of
the four forward differences. Intersect the four tensor-recording sets
with the four difference-recording sets at $j=0$ and with this identity.
Linearity gives
\[
 \left|\lambda^{\sigma}_4\left(y^0_k-\sum_d w_{d,k}\right)\right|
 \le\sum_a H^{\sigma,0}_{a,k}+\sum_dD^{\sigma}_{d,k}.
\]
The detector bound makes the difference fast. The already proved fastness
of each $w_{d,k}$ and finiteness of $I$ now imply that $y^0_k$ is fast.
\end{proof}

Under the displayed fixed-scale root-price inputs, the same proof gives the
more explicit bound
\begin{equation}
 \|y^0_k\|+\|y^-_k\|=O\bigl((1+s_k)^5e^{-s_k/8}\bigr).
 \label{eq:ex-branch-raw-rate}
\end{equation}
Indeed, for the finitely many required tests and insertions, each recording
error is $O((1+s_k)^4e^{-s_k/8})$: the fourth power of the affine clock has
polynomial growth and the point-limit tolerance is $a_k=e^{-s_k}$. In the
sampling argument choose instead
\[
 n_k=\left\lceil\frac{e^{s_k}}{1+s_k}\right\rceil,\qquad
 h_k=a_kn_k.
\]
Then $(1+s_k)^{-1}\le h_k\le(1+s_k)^{-1}+e^{-s_k}$, so $h_k\to0$ and
the endpoint error divided by $h_k$ is
$O((1+s_k)^5e^{-s_k/8})$. The moving detectors and finite trace comparisons
now give \eqref{eq:ex-branch-raw-rate}. In particular, the bound is
$O(e^{-\kappa s_k})$ for every $0<\kappa<1/8$. Its constants depend on the
finite tests; the price and point-limit statements remain inputs.

For the remaining weighted step, let $S_{k,u}$ be the finite-volume
transport whose fixed-$k$ separation limit is $B_k$. With the same
accumulated column logarithm $\ell_k$ as above, the selected calibration
gives a fixed cutoff $k_0$ such that, for every $k\ge k_0$
and every $v\in E_4$, the bound
\[
 \|S_{k,u}v\|\le C\widehat\beta_k e^{\ell_k}\|v\|
\]
on an ultrafilter-large set of $u$. Here $C$ depends only on the fixed
finite-dimensional label coordinates. The closed-ball limit gives the
same bound for $B_kv$, and taking the supremum over $\|v\|\le1$ gives
$\|B_k\|\le C\widehat\beta_k e^{\ell_k}$. This passage fixes $k$ and $v$
before taking the separation limit; it does not require a common
ultrafilter-large set for infinitely many $v$. The accumulated-column estimate gives $\ell_k=o(s_k)$, so $e^{\ell_k}$ is
subexponential. The clock is affine in $s_k$; since
$s\le\epsilon^{-1}e^{\epsilon s}$ for $s\ge0$ and $\epsilon>0$,
$\widehat\beta_k$ and hence $\|B_k\|$ are subexponential as well. The
finite-volume column calibration and fixed-scale transport limit are upstream
inputs. The Crown, Eve and weighted-estimate chapters supply parts of the
selected local prices under their construction hypotheses.

\begin{proposition}[Principal and false trace washout]
\label{prop:ex-branch-washout}
Under the recording and transport estimates,
\eqref{eq:ex-branch-moment-input}, the chart identity and projection
equivariance stated above, for each
fixed $p\in\mathbb N$ and each $b\in\{0,-\}$,
\begin{equation}
 \widehat\beta_{k_i}^{\,p}\,
 \operatorname{tr}(\mathsf P_i t^b_{k_i})\longrightarrow0
 \quad\text{in }\mathcal F_4^0.
 \label{eq:ex-branch-washout}
\end{equation}
In particular the assertion holds for $p=1$. The true branch is treated
separately below.
\end{proposition}
\begin{proof}
First, diagonal symmetrization changes nothing. Let $\mathcal G$ be the
finite signed and real-structure symmetry set, and let $\rho_g$ denote
the corresponding real-linear field action. In each term of the symmetry
average the square of the tensor sign is one and the axis permutation is
bijective. Using projection equivariance gives
\[
 q^b_k:=\operatorname{tr}t^b_k
 =\frac1{|\mathcal G|}\sum_{g\in\mathcal G}\mathfrak c\,\rho_g(y^b_k).
\]
These fixed real-linear maps are bounded in finite dimension. Thus a
constant $C$ independent of $k$ satisfies
$\|q^b_k\|\le C\|y^b_k\|$, and
Proposition~\ref{prop:ex-branch-raw-fast} makes $q^b_k$ fast. It is centered, because each
$t^b_k(a,c)$ is centered and centering is linear.

Choose a fixed linear representative $R:\mathcal F_4^{\mathrm{int}}\to E_4$ with
$\pi R=\mathrm{id}$ and set $r^b_k=Rq^b_k$. It is also fast. For the
fixed integer $p$, subexponential growth of $\widehat\beta^p_kB_k$
therefore gives
\[
 \widehat\beta_k^{\,p}B_kr^b_k\longrightarrow0.
\]
For example, if $\|r^b_k\|\le C e^{-\eta s_k}$, use the subexponential
bound with exponent $\eta/2$ on the combined factor. Composing with the
cofinal sequence $k_i$ preserves the limit. The norms of $\mathsf M_i$
are eventually bounded by \eqref{eq:ex-branch-moment-input}. Hence their
images of the last sequence tend to zero. Since $q^b_k$ is centered and
$\pi Rq^b_k=q^b_k$, the chart identity yields, for all sufficiently large
$i$,
\[
 \mathsf M_i(\widehat\beta_{k_i}^{\,p}B_{k_i}r^b_{k_i})
 =\widehat\beta_{k_i}^{\,p}\mathsf P_iq^b_{k_i}
 =\widehat\beta_{k_i}^{\,p}\operatorname{tr}(\mathsf P_i t^b_{k_i}).
\]
This proves \eqref{eq:ex-branch-washout}. The estimate is for each fixed
$p$; it is not a uniform assertion over a growing exponent.
\end{proof}

With the explicit root-price and point-limit inputs, the quantitative raw
rate gives more: for each fixed $p\in\mathbb N$ and $b\in\{0,-\}$,
\[
 \bigl\|\widehat\beta_{k_i}^{\,p}\operatorname{tr}(\mathsf P_i t^b_{k_i})\bigr\|
 =O\!\left((1+s_{k_i})^{p+6}e^{\ell_{k_i}-s_{k_i}/8}\right).
\]
Indeed, the raw bound contributes $(1+s)^5e^{-s/8}$, while
$\widehat\beta^p\|B\|$ is bounded by a constant times
$(1+s)^{p+1}e^\ell$; the moment maps are eventually bounded. Since
$\ell_k=o(s_k)$, the right side is $O(e^{-\kappa s_{k_i}})$ for every
$0<\kappa<1/8$. The constants can depend on $p$ and $\kappa$.

The true branch instead has the unweighted limiting trace $4\mathcal O$ by
Proposition~\ref{prop:ex-branch-moments} and
\eqref{eq:ex-branch-three-traces}. This limit is nonzero, so the unweighted trace norm is eventually at least
$2\|\mathcal O\|>0$. Since $\widehat\beta_{k_i}\to\infty$, the norm of
the true-branch trace multiplied by $\widehat\beta_{k_i}^{\,p}$ diverges
for each fixed $p\ge1$. These consequences do not construct the physical tensor or close the
quantitative chart inputs.
